% GENERATED FILE. Edit canonical lessons, then run npm run generate:books.
% canonical-source-hash: fnv1a64:5de5526abd628c39
% canonical-lessons: JA-R12-foundation-01, JA-C12-atama, JA-C12-kami, JA-R12-foundation-02, JA-C12-ha-tooth, JA-C12-kata, JA-R12-mixed-scripts, JA-C12-onaka, JA-C12-senaka, JA-R12-doorway, JA-C12-koshi, JA-R12-farewell, JA-R12-repair-01, JA-R12-repair-02, JA-R12-repair-03, JA-R12-repair-04, JA-R12-writing-01, JA-R12-writing-02, JA-R12-body-01, JA-C12-body-map-two

\chapter{Seven More Body Words, Spaced on Purpose}
\label{ch:ja-body-seven-more}
\hlchaptermodality{12}

\begin{chapteropening}
\textbf{By the end of this chapter:} \emph{I can read, write, hear, and say seven more body words while retrieving older script, greeting, farewell, and repair knowledge at expanding intervals.}

Seven fully decodable body labels interleaved with measured R3 and R4 retrieval rather than taught as a single topical list.
\end{chapteropening}

\section[Practice]{Retrieval 1 — mora, politeness, and \ja{し}}

\label{lesson:JA-R12-foundation-01}



\begin{quote}
Say \textbf{\ja{い}} and count its one mora aloud. \emph{Write it:} \textbf{\ja{い}} and \textbf{\ja{し}} from memory

Hear and say the first seven-word body map. \emph{Read it:} the first seven-word body map

\emph{Write it:} the first seven-word body map

Finish by isolating \textbf{\ja{あし}}.
\end{quote}

\subsection*{Your turn}
Hear a polite thank-you and say it. \emph{Read it:} \textbf{\ja{ありがとうございます}}

Then choose it for someone outside your close circle.

\begin{tcolorbox}[breakable,colback=teal!4,colframe=teal!35!black,title={Before you move on}]
Stop after one accurate round.
\end{tcolorbox}
\section[\texorpdfstring{\ja{あたま} (atama)}{あたま (atama)}]{\ja{あたま} — head}

\label{lesson:JA-C12-atama}



\begin{quote}
\emph{Write it:} \textbf{\ja{は}} and \textbf{\ja{日}}

\emph{Read it:} \textbf{\ja{ございます}}

Say \emph{yukkuri}, then retrieve \textbf{\ja{はな}}.
\end{quote}

\subsection*{What to know first}
\begin{quote}
\textbf{\ja{あたま}} — \emph{a-ta-ma} — head
\end{quote}

Every sign is earned. Keep three even morae.

\subsection*{Your turn}
Hear, picture the part, and say \textbf{\ja{あたま}}. \emph{Read it:} \textbf{\ja{あたま}}

\emph{Write it:} \textbf{\ja{あたま}}, with the word hidden

\begin{tcolorbox}[breakable,colback=teal!4,colframe=teal!35!black,title={Before you move on}]
Source: \href{https://words.marugotoweb.jp/}{Japan Foundation Marugoto A1 vocabulary database}.
\end{tcolorbox}
\section[\texorpdfstring{\ja{かみ} (kami)}{かみ (kami)}]{\ja{かみ} — hair}

\label{lesson:JA-C12-kami}



\begin{quote}
\emph{Write it:} \textbf{\ja{あたま}}

\emph{Read it:} \textbf{\ja{はい}}

\emph{Write it:} \textbf{\ja{本}}, then sketch the speech component, then \textbf{\ja{ゆ}}, and finish with \textbf{\ja{かお}}
\end{quote}

\subsection*{What to know first}
\begin{quote}
\textbf{\ja{かみ}} — \emph{ka-mi} — hair
\end{quote}

Two known signs, two morae.

\subsection*{Your turn}
Hear, picture the part, and say \textbf{\ja{かみ}}. \emph{Read it:} \textbf{\ja{かみ}}

\emph{Write it:} \textbf{\ja{かみ}} from memory

\begin{tcolorbox}[breakable,colback=teal!4,colframe=teal!35!black,title={Before you move on}]
Source: \href{https://words.marugotoweb.jp/}{Japan Foundation Marugoto A1 vocabulary database}.
\end{tcolorbox}
\section[Practice]{Retrieval 2 — old signs, recent request}

\label{lesson:JA-R12-foundation-02}



\begin{quote}
\emph{Write it:} \textbf{\ja{え}} and \textbf{\ja{かみ}}, then build the five and mouth components of \textbf{\ja{語}}

Then ask someone to say it with \emph{itte kudasai}. \emph{Write it:} \textbf{\ja{め}}
\end{quote}

\begin{tcolorbox}[breakable,colback=teal!4,colframe=teal!35!black,title={Before you move on}]
Stop after one pass.
\end{tcolorbox}
\section[\texorpdfstring{\ja{は} (ha)}{は (ha)}]{\ja{は} — tooth}

\label{lesson:JA-C12-ha-tooth}



\begin{quote}
Say \textbf{\ja{はい}}. \emph{Write it:} \textbf{\ja{語}}

Recall its Chinese-derived bridge. \emph{Write it:} \textbf{\ja{もうすこし}}

Then hear and say the first seven body words. \emph{Read it:} the first seven body words

\emph{Write it:} the first seven body words
\end{quote}

\subsection*{What to know first}
\begin{quote}
\textbf{\ja{は}} — \emph{ha} — tooth; teeth
\end{quote}

One sign is the complete word.

\subsection*{Your turn}
Hear, picture the part, and say \textbf{\ja{は}}. \emph{Read it:} \textbf{\ja{は}}

\emph{Write it:} \textbf{\ja{は}}

\begin{tcolorbox}[breakable,colback=teal!4,colframe=teal!35!black,title={Before you move on}]
Source: \href{https://words.marugotoweb.jp/}{Japan Foundation Marugoto A1 vocabulary database}.
\end{tcolorbox}
\section[\texorpdfstring{\ja{かた} (kata)}{かた (kata)}]{\ja{かた} — shoulder}

\label{lesson:JA-C12-kata}



\begin{quote}
\emph{Write it:} \textbf{\ja{は}}

Say \textbf{\ja{いいえ}} with its full vowel length. \emph{Read it:} \textbf{\ja{日本語}} with its learned readings

Then close with \emph{wakarimashita}.
\end{quote}

\subsection*{What to know first}
\begin{quote}
\textbf{\ja{かた}} — \emph{ka-ta} — shoulder
\end{quote}

Two known signs, two even morae.

\subsection*{Your turn}
Hear, picture the part, and say \textbf{\ja{かた}}. \emph{Read it:} \textbf{\ja{かた}}

\emph{Write it:} \textbf{\ja{かた}}

\begin{tcolorbox}[breakable,colback=teal!4,colframe=teal!35!black,title={Before you move on}]
Source: \href{https://words.marugotoweb.jp/}{Japan Foundation Marugoto A1 vocabulary database}.
\end{tcolorbox}
\section[Practice]{Retrieval 3 — keep the scripts apart}

\label{lesson:JA-R12-mixed-scripts}



\begin{quote}
\emph{Write it:} hiragana \textbf{\ja{こ}}, then katakana \textbf{\ja{コー}} — the bar holds the vowel

Picture the spot where you are and say \textbf{\ja{ここ}}. Name shoulder, then retrieve head after the longer gap.
\end{quote}

\begin{tcolorbox}[breakable,colback=teal!4,colframe=teal!35!black,title={Before you move on}]
Stop after one pass.
\end{tcolorbox}
\section[\texorpdfstring{\ja{おなか} (onaka)}{おなか (onaka)}]{\ja{おなか} — belly}

\label{lesson:JA-C12-onaka}



\begin{quote}
\emph{Write it:} \textbf{\ja{ん}}, \textbf{\ja{ー}}, and \textbf{\ja{ヒ}}

Ask for slower speech. \emph{Write it:} \textbf{\ja{かみ}}, after its longer gap
\end{quote}

\subsection*{What to know first}
\begin{quote}
\textbf{\ja{おなか}} — \emph{o-na-ka} — belly; stomach area
\end{quote}

Keep three morae.

\subsection*{Your turn}
Hear, picture the part, and say \textbf{\ja{おなか}}. \emph{Read it:} \textbf{\ja{おなか}}

\emph{Write it:} \textbf{\ja{おなか}}

\begin{tcolorbox}[breakable,colback=teal!4,colframe=teal!35!black,title={Before you move on}]
Source: \href{https://words.marugotoweb.jp/}{Japan Foundation Marugoto A1 vocabulary database}.
\end{tcolorbox}
\section[\texorpdfstring{\ja{せなか} (senaka)}{せなか (senaka)}]{\ja{せなか} — back}

\label{lesson:JA-C12-senaka}



\begin{quote}
\emph{Write it:} \textbf{\ja{おなか}} and \textbf{\ja{に}}

\emph{Read it:} \textbf{\ja{コーヒー}} as a katakana loan

\emph{Write it:} small \textbf{\ja{っ}}
\end{quote}

\subsection*{What to know first}
\begin{quote}
\textbf{\ja{せなか}} — \emph{se-na-ka} — back
\end{quote}

Only the opening changes from \textbf{\ja{おなか}}.

\subsection*{Your turn}
Hear, picture the part, and say \textbf{\ja{せなか}}. \emph{Read it:} \textbf{\ja{せなか}}

\emph{Write it:} \textbf{\ja{せなか}}

\begin{tcolorbox}[breakable,colback=teal!4,colframe=teal!35!black,title={Before you move on}]
Source: \href{https://words.marugotoweb.jp/}{Japan Foundation Marugoto A1 vocabulary database}.
\end{tcolorbox}
\section[Practice]{Retrieval 4 — leave the body-word topic}

\label{lesson:JA-R12-doorway}



\begin{quote}
\emph{Write it:} \textbf{\ja{ち}}

Run the old doorway exchange and end with \textbf{\ja{さようなら}}. \emph{Write it:} \textbf{\ja{て}}

Retrieve \textbf{\ja{せなか}}, then \textbf{\ja{は}} after five lessons.
\end{quote}

\begin{tcolorbox}[breakable,colback=teal!4,colframe=teal!35!black,title={Before you move on}]
Stop after one pass.
\end{tcolorbox}
\section[\texorpdfstring{\ja{こし} (koshi)}{こし (koshi)}]{\ja{こし} — waist or lower back}

\label{lesson:JA-C12-koshi}



\begin{quote}
\emph{Write it:} \textbf{\ja{わ}}

Explain why topic \textbf{\ja{は}} can sound \emph{wa}. Hear \emph{sayōnara}. \emph{Write it:} its \textbf{\ja{ら}}, then \textbf{\ja{た}}

Retrieve \textbf{\ja{かた}}.
\end{quote}

\subsection*{What to know first}
\begin{quote}
\textbf{\ja{こし}} — \emph{ko-shi} — waist; lower back
\end{quote}

Two signs, two morae.

\subsection*{Your turn}
Hear, picture the part, and say \textbf{\ja{こし}}. \emph{Read it:} \textbf{\ja{こし}}

\emph{Write it:} \textbf{\ja{こし}}

\begin{tcolorbox}[breakable,colback=teal!4,colframe=teal!35!black,title={Before you move on}]
Source: \href{https://words.marugotoweb.jp/}{Japan Foundation Marugoto A1 vocabulary database}.
\end{tcolorbox}
\section[Practice]{Retrieval 5 — enter, leave, voice}

\label{lesson:JA-R12-farewell}



\begin{quote}
\emph{Read it:} \textbf{\ja{こんにちは}}

Hear and say \textbf{\ja{さようなら}}. \emph{Read it:} \textbf{\ja{さようなら}}

\emph{Write it:} \textbf{\ja{さようなら}}, then \textbf{\ja{だ}}

Retrieve \textbf{\ja{こし}} without looking back.
\end{quote}

\begin{tcolorbox}[breakable,colback=teal!4,colframe=teal!35!black,title={Before you move on}]
Stop after one pass.
\end{tcolorbox}
\section[Practice]{Retrieval 6 — greet, repair, point}

\label{lesson:JA-R12-repair-01}



\begin{quote}
Greet, then say that you do not understand well. \emph{Write it:} \textbf{\ja{て}}

Retrieve \textbf{\ja{おなか}} after five lessons.
\end{quote}

\begin{tcolorbox}[breakable,colback=teal!4,colframe=teal!35!black,title={Before you move on}]
Stop after one pass.
\end{tcolorbox}
\section[Practice]{Retrieval 7 — from \ja{あ} to \ja{せなか}}

\label{lesson:JA-R12-repair-02}



\begin{quote}
\emph{Write it:} \textbf{\ja{あ}} and \textbf{\ja{ど}}

Say \emph{mō} on two morae. \emph{Write it:} \textbf{\ja{みみ}}

Then retrieve \textbf{\ja{せなか}} after five lessons.
\end{quote}

\begin{tcolorbox}[breakable,colback=teal!4,colframe=teal!35!black,title={Before you move on}]
Stop after one pass.
\end{tcolorbox}
\section[Practice]{Retrieval 8 — \ja{り}, \ja{も}, \ja{いちど}, \ja{くち}}

\label{lesson:JA-R12-repair-03}



\begin{quote}
\emph{Write it:} \textbf{\ja{り}} and \textbf{\ja{も}}

Say \textbf{\ja{いちど}}. \emph{Write it:} \textbf{\ja{いちど}}

Then retrieve \textbf{\ja{くち}}.
\end{quote}

\begin{tcolorbox}[breakable,colback=teal!4,colframe=teal!35!black,title={Before you move on}]
Stop after one pass.
\end{tcolorbox}
\section[Practice]{Retrieval 9 — request, leg, waist}

\label{lesson:JA-R12-repair-04}



\begin{quote}
\emph{Write it:} \textbf{\ja{か}}

Say \emph{onegaishimasu}, then the complete repetition request. \emph{Write it:} \textbf{\ja{あし}}

Then retrieve \textbf{\ja{こし}} after five lessons.
\end{quote}

\begin{tcolorbox}[breakable,colback=teal!4,colframe=teal!35!black,title={Before you move on}]
Stop after one pass.
\end{tcolorbox}
\section[Practice]{Retrieval 10 — mark, signs, nose}

\label{lesson:JA-R12-writing-01}



\begin{quote}
\emph{Write it:} a known base with dakuten added, then \textbf{\ja{せ}} and \textbf{\ja{ね}}, then \textbf{\ja{はな}} from memory
\end{quote}

\begin{tcolorbox}[breakable,colback=teal!4,colframe=teal!35!black,title={Before you move on}]
Stop after one pass.
\end{tcolorbox}
\section[Practice]{Retrieval 11 — \ja{と}, \ja{すこし}, \ja{かお}}

\label{lesson:JA-R12-writing-02}



\begin{quote}
\emph{Write it:} \textbf{\ja{と}}, then \textbf{\ja{すこし}}, and finish with \textbf{\ja{かお}}

Say \emph{ka-o} on two morae.
\end{quote}

\begin{tcolorbox}[breakable,colback=teal!4,colframe=teal!35!black,title={Before you move on}]
Stop after one pass.
\end{tcolorbox}
\section[Practice]{Retrieval 12 — \ja{う} and \ja{め}}

\label{lesson:JA-R12-body-01}



\begin{quote}
\emph{Write it:} \textbf{\ja{う}}

Then hear and say \textbf{\ja{め}}. \emph{Read it:} \textbf{\ja{め}}

\emph{Write it:} \textbf{\ja{め}}
\end{quote}

\begin{tcolorbox}[breakable,colback=teal!4,colframe=teal!35!black,title={Before you move on}]
Stop after one pass.
\end{tcolorbox}
\section[\texorpdfstring{\ja{あたま・かみ・は} … (atama, kami, ha …)}{あたま・かみ・は … (atama, kami, ha …)}]{Seven more body words, spaced on purpose}

\label{lesson:JA-C12-body-map-two}



\begin{quote}
\emph{Read it:} \textbf{\ja{ありがとう}}

Then hear and say any three words from the first seven-word body map. \emph{Read it:} the same three words

\emph{Write it:} the same three words
\end{quote}

\subsection*{What to know first}
\begin{quote}
\textbf{\ja{あたま・かみ・は・かた・おなか・せなか・こし}}
\end{quote}

Every sign is earned, and every new word has already had short and medium retrieval scheduled in this chapter.

\subsection*{Your turn}
Listen and give each English meaning; say all seven. \emph{Read it:} all seven

\emph{Write it:} any four from memory

\begin{tcolorbox}[breakable,colback=teal!4,colframe=teal!35!black,title={Before you move on}]
Stop after one accurate seven-word map.

Source: \href{https://words.marugotoweb.jp/}{Japan Foundation Marugoto A1 vocabulary database}.
\end{tcolorbox}
